\section{The confined isoperimetric profile}
\label{sec:profile}

Confinement gives perimeter minimizers at prescribed volume. At volumes
where the confined profile is differentiable, its derivative will determine
the mean curvature on the free boundary and bound it from above on the
contact set. Both portions contribute to full ambient perimeter. These
curvature bounds will enter the hypersurface Sobolev inequality once the
constant test has been justified in Section~\ref{sec:comparison}.

Throughout this section, $M^n$ is a Cartan--Hadamard manifold; only
$\Sec_g\le0$ is needed. Fix an open geodesic ball
$B=B_{R_B}(o)$ of finite radius. Its closure is compact, and its
boundary is a smooth strictly convex hypersurface. For a measurable set
$E$, write $|E|=\vol_g(E)$ and
\[
 P(E;U)=|D\chi_E|(U),\qquad P(E)=P(E;M),
\]
where $U\subset M$ is open and $|D\chi_E|$ is the perimeter measure.
In particular, $P(E)$ includes any part of the boundary lying on
$\partial B$.  Define
\begin{equation}
 \label{eq:confined-profile}
 I(v)=\inf\bigl\{P(E):\ |E\setminus B|=0,\ |E|=v\bigr\},
 \qquad 0<v<|B|.
\end{equation}
All statements about measurable sets in this definition are understood
up to sets of ambient volume zero.
\subsection{Compactness and variation of finite-perimeter sets}

For an integrable function $u$, write $u\in BV(M)$ when its
distributional gradient $Du$ is a finite vector-valued measure, and
write $|Du|$ for its total variation. For $u=\chi_E$, this agrees
with the perimeter measure already defined.

Existence of confined minimizers will follow from compactness, and smooth
flows will change their volume while controlling full ambient perimeter.
We recall these facts in the precise forms used below. In smooth coordinate charts,
De Giorgi's reduced-boundary representation and Gauss--Green formula
\cite[Synopsis, pp.~119--120]{Maggi2012} give the Riemannian identities
\[
 D\chi_E=-N_E\,\mathcal H_g^{n-1}\lfloor\partial^*E,
 \qquad
 \int_E\divg X\,d\vol_g
 =\int_{\partial^*E}\langle X,N_E\rangle\,d\mathcal H_g^{n-1}.
\]
Here $\partial^*E$ is countably rectifiable, $N_E$ is its outward
measure-theoretic unit normal, and its tangent plane is $N_E^\perp$
almost everywhere. The formulas apply to every smooth compactly
supported field $X$.
To see the metric conversion, write $g$ for the metric matrix,
$\rho=\sqrt{\det g}$ and $\nu$ for the Euclidean outward normal.
If $a=(\nu^Tg^{-1}\nu)^{1/2}$, then
\[
 N_E=g^{-1}\nu/a,\qquad
 d\mathcal H_g^{n-1}=\rho a\,d\mathcal H_e^{n-1},\qquad
 \divg_gX=\rho^{-1}\divg(\rho X).
\]
The first two identities follow from orthogonality and the tangent
Gram determinant; substitution in Euclidean Gauss--Green gives the
displayed Riemannian formula. Smooth positivity of the metric and
density compares the BV norms on compact subcharts. Chart partitions
then give the invariant measure identity.

\begin{lemma}[Compactness and lower semicontinuity]\label{lem:bv-compactness}
Let $K$ be a compact subset of a smooth Riemannian manifold without
boundary. If functions $u_j$ vanish outside $K$ almost everywhere and
\[
 \sup_j\bigl(\|u_j\|_{L^1}+|Du_j|(M)\bigr)<\infty,
\]
then a subsequence converges in $L^1(M)$. Total variation is lower
semicontinuous under this convergence. If the $u_j$ are characteristic
functions, their limit is a characteristic function.
\end{lemma}

\begin{proof}
Choose a finite smooth partition on a neighborhood of $K$, with each
partition function compactly supported in a coordinate chart.
The smooth multiplier formula follows by inserting a multiplier in
the distributional test field. Smooth positivity of the metric and
volume density on the compact chart supports therefore bounds the
Euclidean BV norms of the localized functions. Extend each of them
by zero within its chart.

For a Euclidean BV function $v$, convolution satisfies
$D(v*\rho_\varepsilon)=Dv*\rho_\varepsilon$ and
$\|D(v*\rho_\varepsilon)\|_1\le |Dv|(\R^n)$.
Apply the fundamental theorem of calculus to the convolution and let
its scale tend to zero in $L^1$. This gives
\[
 \|v(\cdot+h)-v\|_1\le |h|\,|Dv|(\R^n),\qquad
 \|v*\rho_\varepsilon-v\|_1
 \le\varepsilon C_\rho |Dv|(\R^n).
\]
For fixed $\varepsilon$, the convolutions have common compact support
and uniformly bounded sup norms and first derivatives, by the uniform
$L^1$ bound on $v$. Approximation on a finite mesh makes this family
precompact in $L^1$. Taking scales to zero and a diagonal subsequence
gives compactness in each chart, hence for the original finite sum.
Total variation is a supremum of continuous divergence integrals, so
is lower semicontinuous. A further almost-everywhere convergent
subsequence of characteristic functions has values in $\{0,1\}$.
\end{proof}

\begin{lemma}[Perimeter under smooth flows]\label{lem:bv-flow}
Let $F:M\to M$ be an orientation-preserving smooth diffeomorphism,
and let $E$ have finite perimeter. If the right side below is finite,
then
\begin{equation}\label{eq:bv-flow-jacobian}
 P(FE)=\int_{\partial^*E}
 J_{n-1}\bigl(dF_x|_{N_E(x)^\perp}\bigr)\,d\mathcal H_g^{n-1}(x).
\end{equation}
All Jacobians and adjoints use the source and target Riemannian
metrics. For a smooth flow $F_t$ with velocity $Z$ and a set whose
perimeter support is compact,
\begin{equation}\label{eq:bv-flow-first-variation}
 \left.\frac d{dt}\right|_{t=0}P(F_tE)
 =\int_{\partial^*E}\divg_{N_E^\perp}Z\,d\mathcal H_g^{n-1}.
\end{equation}
For a fixed finite family of flows and perimeter supports contained
in a fixed compact set, there are common $C,\tau>0$ such that
$|P(F_tE)-P(E)|\le C|t|P(E)$ for $|t|<\tau$.
\end{lemma}

\begin{proof}
Write $J_F$ for the positive ambient volume Jacobian. The Piola
pullback of a target test field $X$ is
\[
 Q_X(x)=J_F(x)(dF_x)^{-1}X(Fx).
\]
Ordinary smooth change of variables, tested against smooth scalar
functions, gives
$\divg Q_X=J_F(\divg X)\circ F$. Gauss--Green for $E$ implies
\[
 \int_{FE}\divg X\,d\vol_g
 =\int_{\partial^*E}
 \langle X(Fx),J_F(x)(dF_x)^{-*}N_E(x)\rangle_g
 \,d\mathcal H_g^{n-1}(x).
\]
Thus the outward vector measure $-D\chi_{FE}$ is the pushforward of
the vector density $a(x)=J_F(x)(dF_x)^{-*}N_E(x)$ against the
perimeter measure of $E$. Since $F$ is injective and $a\ne0$, its
polar direction at $y=F(x)$ is $a(x)/|a(x)|$. Its total variation is
therefore the pushforward of $|a|\,d|D\chi_E|$, with no cancellation.
Linear algebra in orthonormal bases gives
$|a|=J_{n-1}(dF|_{N_E^\perp})$, proving
\eqref{eq:bv-flow-jacobian}. Differentiation of the Gram determinant
at $t=0$ gives tangential divergence. The Jacobians and their time
derivatives are uniformly bounded over all tangent planes above the
fixed compact set. This justifies differentiation under the integral
and proves the final estimate.
\end{proof}

The compactness statement supplies minimizers and compact families of
competitors. The flow formula counts contact perimeter as well as free
perimeter, both for volume adjustment and for one-sided variations.


\subsection{Existence, regularity, and variation of the profile}

\begin{proposition}[Confined minimizers and their regularity]
 \label{prop:confined-regularity}
For each $v\in(0,|B|)$ there is a minimizer $E$ in
\eqref{eq:confined-profile}.  It has a regular representative for which
$\Gamma=\partial E$ is the support of its perimeter measure and is
contained in $\overline B$.  There is a compact set $S\subset B$ such
that
\begin{enumerate}
 \item $\dim_{\mathcal H}S\le n-8$, with $S=\varnothing$ if $n<8$;
 \item $\Sigma=\Gamma\setminus S$ is a boundaryless
 $C^{1,1}_{\mathrm{loc}}$ hypersurface, smooth on $\Sigma\cap B$;
 \item $\Gamma$ is $C^{1,1}$ in an ambient neighborhood of
 $\partial B$, and $S$ is disjoint from a collar of $\partial B$;
 \item with $d\sigma$ denoting area on $\Sigma$, for every Borel set
 $A\subset M$ one has
 \begin{equation}
  \label{eq:perimeter-area}
  |D\chi_E|(A)=\int_{A\cap\Sigma}d\sigma,
  \qquad P(E)=\int_\Sigma d\sigma=I(v).
 \end{equation}
\end{enumerate}
The unit normal on $\Sigma$ is the outward normal of $E$.
\end{proposition}

\begin{proof}
Concentric balls provide finite-perimeter competitors at every volume
in $(0,|B|)$.  A minimizing sequence has uniformly bounded perimeter
and is supported, up to null sets, in the compact set $\overline B$.
Lemma~\ref{lem:bv-compactness} gives an $L^1(M)$ limit
$\chi_E$ of the same volume with $P(E)\le I(v)$.  The limit vanishes
almost everywhere outside $\overline B$; since $\partial B$ has
ambient volume zero, it is admissible in \eqref{eq:confined-profile}.
Thus it attains the infimum.

For the regularity conclusions, choose a minimizing region supplied by
Ghomi--Spruck in the cited arXiv version
\cite[Lemma~7.2(i),(ii), p.~35]{GhomiSpruck2022}, and denote it by $E$.
Set $K=\supp|D\chi_E|$. These unconditional parts of the Lemma give interior
smoothness outside the stated singular set and
$C^{1,1}$ regularity in a neighborhood of the confining sphere in a
Cartan--Hadamard manifold; see also
Morgan~\cite[Section~2 and Corollary~3.8]{Morgan2003} for the interior theory.
For the boundary part, Ghomi and Spruck use the Euclidean obstacle
regularity of Stredulinsky and Ziemer \cite{StredulinskyZiemer1997}
and its local graph extension to the Riemannian setting.
The graph and singular-set description gives $|K|=0$.
Off $K$, $\chi_E$ is locally constant almost everywhere.
Take $E$ to be the union of the components of $M\setminus K$ on which
$\chi_E=1$ almost everywhere.  This changes only a null set.
Every ball centered on $K$ has positive perimeter, hence contains both
phases in positive volume, so $\partial E=K=\Gamma$.
At regular points the perimeter measure is the induced area measure.
The singular set has zero $m$-dimensional Hausdorff measure, giving
\eqref{eq:perimeter-area}.

The free portion in the $C^{1,1}$ collar is smooth by interior
regularity for the constant-mean-curvature equation.  Hence the actual
interior singular set avoids this collar.  It is closed away from the
collar and contained in the compact set $\Gamma$, so it is a compact
subset of $B$.  Finally, the $C^{1,1}$ regularity is that of the whole
hypersurface near $\partial B$; the transition between its free and
contact portions does not create a boundary of $\Sigma$.
\end{proof}

We henceforth use this representative whenever discussing the boundary
of a minimizer.  Notice that
\begin{equation}
 \label{eq:positive-free-perimeter}
 P(E;B)>0 \qquad (0<|E|<|B|).
\end{equation}
Indeed, if $D\chi_E$ vanished in the connected ball $B$, then
$\chi_E$ would be constant almost everywhere there, contrary to the
volume condition.  In particular $I(v)>0$, and, by
\eqref{eq:perimeter-area}, a minimizing boundary has a nonempty smooth
free patch.

\begin{lemma}[Local Lipschitz continuity]
 \label{lem:profile-lipschitz}
The profile $I$ is positive and locally Lipschitz on $(0,|B|)$.
\end{lemma}

\begin{proof}
Positivity was just established.  Fix a compact interval
$[v_-,v_+]\subset(0,|B|)$.  The radius of a concentric ball varies
continuously with its volume, and its boundary area is continuous in
its radius.  Consequently there is $C_0<\infty$ such that
$I(v)\le C_0$ for all $v\in[v_-,v_+]$.

Consider the family
\[
 \mathcal K=\bigl\{\chi_F:\ |F\setminus B|=0,
       \ v_-\le |F|\le v_+,\ P(F)\le C_0\bigr\}.
\]
BV compactness on a compact neighborhood of $\overline B$, together
with lower semicontinuity and $L^1$ continuity of volume, shows that
$\mathcal K$ is compact in $L^1(M)$.  For each $F\in\mathcal K$ there
is a smooth vector field $Z_F$ compactly supported in $B$ for which
\[
 A_{Z_F}(F):=\int_F\divg Z_F\,d\vol_g\ne0.
\]
Otherwise $D\chi_F$ would vanish in $B$, contradicting
$v_-\le |F|\le v_+$.  Reverse the sign of $Z_F$ if necessary so that
$A_{Z_F}(F)>0$.  For each fixed $Z$, the functional $A_Z$ is
$L^1$-continuous.  A finite subcover of $\mathcal K$ therefore gives
smooth fields $Z_1,\ldots,Z_k$, compactly supported in $B$, and a
number $\delta>0$ such that
\begin{equation}
 \label{eq:uniform-volume-speed}
 \max_{1\le j\le k} A_{Z_j}(F)\ge 2\delta
 \qquad\text{for every }F\in\mathcal K.
\end{equation}

Let $\Psi_j^t$ be the flow of $Z_j$.  These flows preserve $B$ in
both time directions.  If $J_j(t,x)$ is the ambient volume Jacobian,
then
\[
 G_j(t,F):=|\Psi_j^t(F)|=\int_F J_j(t,x)\,d\vol_g(x),
 \qquad \partial_tG_j(0,F)=A_{Z_j}(F).
\]
Smoothness of the finitely many flows gives uniform bounds on
$\partial_t^2J_j$ on a fixed compact set for small $|t|$.
Since $|F|\le |B|$, there is a common $\tau>0$ such that
\[
 |\partial_tG_j(t,F)-A_{Z_j}(F)|\le\delta
 \quad (|t|\le\tau,\ F\in\mathcal K,\ 1\le j\le k).
\]
For an index satisfying \eqref{eq:uniform-volume-speed}, the function
$t\mapsto G_j(t,F)$ thus has derivative at least $\delta$ on
$[-\tau,\tau]$.  It follows that every $w$ with
$|w-|F||<\delta\tau$ is attained for some $t$ satisfying
\begin{equation}
 \label{eq:volume-adjustment-time}
 |\Psi_j^t(F)|=w,\qquad |t|\le \delta^{-1}|w-|F||.
\end{equation}

By Lemma~\ref{lem:bv-flow}, perimeter transforms by the tangential
Jacobian of a smooth diffeomorphism on the reduced boundary.  Applied to these finitely many
flows, this gives
\begin{align*}
 P(\Psi_j^t(F))
 &=\int_{\partial^*F}
 J_m\bigl(D\Psi_j^t(x)|_{T_x\partial^*F}\bigr)
 \,d\mathcal H_g^m(x),\\
 |P(\Psi_j^t(F))-P(F)|&\le C|t|P(F)
 \qquad (|t|\le\tau).
\end{align*}
Here $C$ is uniform over the finitely many flows and all $m$-planes
above the compact set $\overline B$.
Equation~\eqref{eq:volume-adjustment-time} and the bound
$P(F)\le C_0$ show that changing the volume by $\Delta v$ costs at
most $C C_0\delta^{-1}|\Delta v|$ in perimeter.

Apply this construction first to a minimizer at $v$ and then to a
minimizer at $w$, where $v,w\in[v_-,v_+]$ and $|v-w|<\delta\tau$.
Both are in $\mathcal K$, and the resulting competitors give
\[
 |I(v)-I(w)|\le C C_0\delta^{-1}|v-w|.
\]
This proves local Lipschitz continuity.  All constants here may depend
on the fixed ball and interval; no uniform regularity of the
minimizers is required.
\end{proof}

The profile can now be differentiated at almost every volume. The next
lemma identifies its derivative on free boundary patches and obtains the
needed one-sided bound at contact. Figure~\ref{fig:confined-boundary}
shows why both parts must enter the same ambient perimeter; in the
figure, $\rho=d(o,\cdot)$.

\begin{figure}[htbp]
\centering
\begin{tikzpicture}[scale=.95,>=Latex]
 \draw[gray!65,dashed] (0,0) circle (2.2);
 \fill[blue!7] (75:2.2) arc (75:105:2.2)
   .. controls (-1.65,1.84) and (-1.6,.5) .. (-1.05,-.45)
   .. controls (-.55,-1.35) and (.75,-1.4) .. (1.2,-.55)
   .. controls (1.7,.4) and (1.65,1.84) .. (75:2.2) -- cycle;
 \draw[thick,blue!70!black] (105:2.2)
   .. controls (-1.65,1.84) and (-1.6,.5) .. (-1.05,-.45)
   .. controls (-.55,-1.35) and (.75,-1.4) .. (1.2,-.55)
   .. controls (1.7,.4) and (1.65,1.84) .. (75:2.2);
 \draw[very thick,orange!80!black] (75:2.2) arc (75:105:2.2);
 \node at (0,.3) {$E$};
 \node[gray!75!black] at (2.55,-1.8) {$\partial B$};
 \draw[->,gray!65] (2.2,-1.75)--(1.66,-1.45);
 \draw[->,thick] (0,2.2)--(0,2.85) node[right] {$N=\nabla\rho$};
 \node[orange!80!black,anchor=west] at (1.05,2.45) {$C=\Gamma\cap\partial B$};
 \draw[->,orange!80!black] (1,2.4)--(.42,2.17);
 \node[blue!70!black,anchor=west] at (2.3,.45) {free part $\Sigma\cap B$};
 \draw[->,blue!70!black] (2.3,.32)--(1.47,.15);
\end{tikzpicture}
\caption{The confined problem minimizes full ambient perimeter, including
contact with the confining sphere. The occupied phase fixes the outward
normal. Free variations give $mh=I'(v)$; inward variations give
$mh\le I'(v)$ almost everywhere on the contact set. The planar drawing is schematic:
the proof does not assume that the contact set contains an open patch.}
\label{fig:confined-boundary}
\end{figure}


\begin{lemma}[Mean curvature at differentiability volumes]
 \label{lem:profile-curvature}
Let $v\in(0,|B|)$ be a differentiability point of $I$, and let $E$ be
a minimizer as in Proposition~\ref{prop:confined-regularity}.  With
$mh=\divg_\Sigma N$ and $h_0=I'(v)/m$, one has
\[
 I'(v)>0,\qquad h=h_0\ \text{on }\Sigma\cap B,
 \qquad 0\le h\le h_0\quad\text{almost everywhere on }\Sigma.
\]
In particular, $h^2\le h_0^2$ almost everywhere on $\Sigma$.
\end{lemma}

\begin{proof}
For a smooth flow $\Psi^t$ with $E_t=\Psi^t(E)$, set
$P(t)=P(E_t)$ and $V(t)=|E_t|$.  Whenever the flow is admissible,
\begin{equation}
 \label{eq:profile-variation-comparison}
 P(t)\ge I(V(t)),\qquad P(0)=I(v),\qquad V(0)=v.
\end{equation}
For a field supported on a free regular patch, admissibility holds
for both signs of $t$.  Differentiating at the minimum in
\eqref{eq:profile-variation-comparison} gives
$P'(0)=I'(v)V'(0)$.  The first variation formulas are
\[
 V'(0)=\int_\Sigma\langle Z,N\rangle\,d\sigma,
 \qquad
 P'(0)=m\int_\Sigma h\langle Z,N\rangle\,d\sigma.
\]
Arbitrary smooth compactly supported normal speeds on the patch
therefore give $mh=I'(v)$.  This holds on every free regular patch,
so $h=h_0$ throughout $\Sigma\cap B$.  Such patches exist by
\eqref{eq:positive-free-perimeter}.

To determine the sign of $h_0$, put $\rho(x)=d(o,x)$ and
$F(x)=\rho(x)^2/2$.  The function $F$ is smooth on $M$, including
at $o$.  Hessian comparison gives
\[
 \Hess F=d\rho\otimes d\rho+\rho\Hess\rho
 \ge d\rho\otimes d\rho
       +(g-d\rho\otimes d\rho)
 =g,
\]
with the inequality at $o$ understood by smooth extension.
The flow of $-\nabla F$ is
\[
 \Psi^t(x)=\exp_o\bigl(e^{-t}\exp_o^{-1}(x)\bigr).
\]
It sends $B$ into itself for $t\ge0$.  The volume and perimeter
Jacobian formulas give
\begin{align}
 \label{eq:radial-volume-derivative}
 V'(0)&=-\int_E\Delta F\,d\vol_g\le -nv<0,\\
 \label{eq:radial-perimeter-derivative}
 P'(0)&=-\int_{\partial^*E}
          \tr_{T_x\partial^*E}(\Hess F)\,d\mathcal H_g^m(x)
          \le -mP(E)<0.
\end{align}
These formulas apply to the reduced boundary as a rectifiable set,
so they include its contact portion and do not require smoothness
at singular points.  The integrands and their flow derivatives are
bounded on a fixed compact neighborhood of $\overline B$.
The one-sided derivative of
\eqref{eq:profile-variation-comparison} now yields
\[
 P'(0)-I'(v)V'(0)\ge0.
\]
Since $V'(0)<0$, division reverses the inequality and gives
\begin{equation}
 \label{eq:profile-derivative-positive}
 I'(v)\ge\frac{P'(0)}{V'(0)}>0.
\end{equation}

It remains to control the contact set $C=\Gamma\cap\partial B$.
At every point of $C$, the $C^1$ boundary of $E$ is tangent to the
sphere and has the same outward normal $N=\nabla\rho$: the region
lies on the inner side of the sphere.  At almost every such point
$x$, its local $C^{1,1}$ graph has a second-order expansion.  In
normal coordinates at $x$, with the common outward normal pointing
upward, write the boundary of $E$ and the sphere as graphs $u$ and
$\beta$.  They have equal values and first derivatives at $x$, and
$u\le\beta$.  Thus $D^2u(x)\le D^2\beta(x)$.  In these coordinates their
second fundamental forms for the outward normal are $-D^2u(x)$ and
$-D^2\beta(x)$, respectively. Hessian comparison on the sphere therefore gives
\begin{equation}
 \label{eq:contact-curvature-lower}
 \mathrm{II}_{\Gamma}(x)\ge\mathrm{II}_{\partial B}(x),
 \qquad h(x)\ge R_B^{-1}>0
 \quad\text{for almost every }x\in C.
\end{equation}

Choose an open collar $U$ of $\partial B$ whose closure avoids $S$
and $o$ and in which $\Gamma$ is $C^{1,1}$.  For arbitrary
$\varphi\in C_c^\infty(U)$ with $\varphi\ge0$, the smooth field
$Z=-\varphi\nabla\rho$, extended by zero, has an admissible flow for
$t\ge0$, because $\rho$ is nonincreasing along its trajectories.
The one-sided variation inequality gives
\begin{equation}
 \label{eq:contact-variation}
 m\int_\Sigma(h-h_0)\langle Z,N\rangle\,d\sigma\ge0.
\end{equation}
For completeness, the perimeter first variation used here is obtained
by integrating
$\divg_\Sigma Z=\divg_\Sigma Z^{\mathsf T}
 +mh\langle Z,N\rangle$ on the whole $C^{1,1}$ hypersurface in
the collar.  The tangential divergence integrates to zero, since
its support is compact in $\Sigma$.  Thus there is no additional
term on the transition between the free and contact portions.

The integrand in \eqref{eq:contact-variation} vanishes on the free
portion, where $h=h_0$, while $\langle Z,N\rangle=-\varphi$ on $C$.
Consequently
\[
 \int_C(h-h_0)\varphi\,d\sigma\le0
 \qquad (\varphi\in C_c^\infty(U),\ \varphi\ge0).
\]
Nonnegative smooth ambient tests determine the sign of a Radon
measure, so $h\le h_0$ almost everywhere on $C$, irrespective of
whether $C$ has relative interior.  Combining this with
\eqref{eq:contact-curvature-lower}, $h=h_0$ on the free portion,
and \eqref{eq:profile-derivative-positive}, gives
\begin{equation}
 \label{eq:profile-curvature-bound}
 0\le h\le h_0,\qquad h^2\le h_0^2
 \quad\text{almost everywhere on }\Sigma.
\end{equation}
\end{proof}
